\begin{table}[ht]
\centering
\caption{Positioning of this paper relative to recent empirical work on the post-ChatGPT shock to programming Q\&A platforms. A check mark indicates that the cited study reports a substantive empirical result on the corresponding margin. ``Task-level DDD'' refers to estimation via a triple difference exploiting variation in task-level LLM substitutability; ``support-post funnel'' refers to tracing the effect through support stages: answered, accepted, accepted and not closed, accepted with non-negative score, and funnel-qualified (all of these gates at once).}
\label{tab:literature_positioning}
\footnotesize
\setlength{\tabcolsep}{3.5pt}
\begin{adjustbox}{max width=\textwidth}
\begin{tabular}{lccccccc}
\toprule
Study & \makecell{Quantity\\decline} & \makecell{Quality /\\composition} & \makecell{Users /\\retention} & Answers & \makecell{Novelty /\\readability} & \makecell{Task-level\\DDD} & \makecell{Support-post\\funnel} \\
\midrule
\citet{delriochanona2024large} & \checkmark & \checkmark & & & & & \\
\citet{burtch2024chatgpt} & \checkmark & \checkmark & \checkmark & & & & \\
\citet{xue2026chatgpt} & \checkmark & \checkmark & \checkmark & \checkmark & & & \\
\citet{borwankar2023unraveling} & \checkmark & \checkmark & \checkmark & \checkmark & & & \\
\citet{shanqiu2025isr} & \checkmark & \checkmark & & \checkmark & & & \\
\citet{quinngutt2024heterogeneous} & \checkmark & \checkmark & \checkmark & & \checkmark & & \\
\textbf{This paper} & \checkmark & \checkmark & & & & \checkmark & \checkmark \\
\bottomrule
\end{tabular}
\end{adjustbox}
\end{table}
